% ============================================================================
% ARBEITSBLATT DS 2 - Gruppe B: Wortbildung (S. 104)
% ============================================================================

\documentclass[11pt,a4paper]{article}

\usepackage[utf8]{inputenc}
\usepackage[T1]{fontenc}
\usepackage[ngerman]{babel}
\usepackage[left=2cm,right=2cm,top=1.8cm,bottom=1.5cm]{geometry}
\usepackage{helvet}
\usepackage[table]{xcolor}
\usepackage{tabularx}
\usepackage{array}
\usepackage{enumitem}
\usepackage{tcolorbox}
\usepackage{fancyhdr}
\usepackage{lastpage}

\renewcommand{\familydefault}{\sfdefault}

\definecolor{spanisch}{HTML}{c41e3a}
\definecolor{grau}{HTML}{666666}
\definecolor{hellgrau}{HTML}{f5f5f5}
\definecolor{gruppeB}{HTML}{1565c0}

\pagestyle{fancy}
\fancyhf{}
\fancyhead[L]{\small\textcolor{grau}{Spanisch Spätbeginner -- Gruppe B}}
\fancyhead[R]{\small\textcolor{grau}{AB-02-B}}
\fancyfoot[C]{\small\thepage/\pageref{LastPage}}
\renewcommand{\headrulewidth}{0.5pt}

\newcommand{\luecke}[1]{\underline{\hspace{#1}}}
\newcommand{\punktebox}[1]{\hfill\fbox{\small \_\_/#1}}
\newcommand{\es}[1]{\textit{#1}}
\newcommand{\pfeil}{$\rightarrow$}

\newtcolorbox{merkkasten}[1][]{
    colback=white,
    colframe=gruppeB,
    fonttitle=\bfseries\color{gruppeB},
    title=#1,
    boxrule=1.5pt,
    arc=0pt,
    left=5pt, right=5pt, top=5pt, bottom=5pt
}

\newtcolorbox{buchverweis}{
    colback=hellgrau,
    colframe=grau,
    boxrule=0.5pt,
    arc=0pt,
    left=5pt, right=5pt, top=3pt, bottom=3pt
}

\newcounter{aufgabe}
\newcommand{\aufgabe}[2]{%
    \stepcounter{aufgabe}%
    \vspace{0.8em}%
    \noindent\textbf{\theaufgabe.} #1 \punktebox{#2}%
    \vspace{0.3em}%
}

\begin{document}

% ============================================================================
% KOPFBEREICH
% ============================================================================
\noindent
\begin{tabularx}{\textwidth}{@{}Xr@{}}
    {\Large\textbf{\textcolor{gruppeB}{Formación de palabras}}} & Name: \luecke{5cm} \\[0.3em]
    {\small DS 2 -- Wortbildung (S. 104)} & Datum: \luecke{3cm} \\
\end{tabularx}

\vspace{0.3em}
\hrule
\vspace{0.8em}

% ============================================================================
% MERKKASTEN 1: Suffixe
% ============================================================================
\begin{merkkasten}[Kasten 1: Wortbildung mit Suffixen]
\vspace{0.2em}

\begin{tabularx}{\textwidth}{@{}|p{2.5cm}|X|X|@{}}
\hline
\rowcolor{hellgrau}
\textbf{Suffix} & \textbf{Bildungsregel} & \textbf{Beispiel} \\
\hline
\textbf{-ción} & Verb \pfeil{} Substantiv & comunicar \pfeil{} comunicación \\[0.3em]
\hline
\textbf{-(i)dad} & Adjektiv \pfeil{} Substantiv & responsable \pfeil{} responsabilidad \\[0.3em]
\hline
\textbf{-ar} & Substantiv \pfeil{} Verb & grupo \pfeil{} agrupar \\[0.3em]
\hline
\end{tabularx}

\vspace{0.3em}
\textbf{Tipp:} Ähnlich wie im Deutschen/Englischen: -tion, -ity, -ieren
\vspace{0.2em}
\end{merkkasten}

\vspace{1em}

% ============================================================================
% AUFGABEN
% ============================================================================

\aufgabe{Bilde Substantive auf \es{-ción}. (Buch S. 104, Aufg. 4a)}{8}

\begin{tabularx}{\textwidth}{@{}|X|X||X|X|@{}}
\hline
\rowcolor{hellgrau}
\textbf{Verb} & \textbf{Substantiv} & \textbf{Verb} & \textbf{Substantiv} \\
\hline
informar & \luecke{3cm} & organizar & \luecke{3cm} \\[0.5em]
\hline
comunicar & \luecke{3cm} & presentar & \luecke{3cm} \\[0.5em]
\hline
formar & \luecke{3cm} & evaluar & \luecke{3cm} \\[0.5em]
\hline
motivar & \luecke{3cm} & colaborar & \luecke{3cm} \\[0.5em]
\hline
\end{tabularx}

\vspace{0.8em}

\aufgabe{Bilde Substantive auf \es{-(i)dad}. (Buch S. 104, Aufg. 4b)}{6}

\begin{tabularx}{\textwidth}{@{}|X|X||X|X|@{}}
\hline
\rowcolor{hellgrau}
\textbf{Adjektiv} & \textbf{Substantiv} & \textbf{Adjektiv} & \textbf{Substantiv} \\
\hline
responsable & \luecke{3cm} & capaz & \luecke{3cm} \\[0.5em]
\hline
posible & \luecke{3cm} & creativo & \luecke{3cm} \\[0.5em]
\hline
flexible & \luecke{3cm} & productivo & \luecke{3cm} \\[0.5em]
\hline
\end{tabularx}

\vspace{0.8em}

\aufgabe{Bilde Verben auf \es{-ar}. (Buch S. 104, Aufg. 4c)}{4}

\begin{tabularx}{\textwidth}{@{}|X|X||X|X|@{}}
\hline
\rowcolor{hellgrau}
\textbf{Substantiv} & \textbf{Verb} & \textbf{Substantiv} & \textbf{Verb} \\
\hline
el grupo & \luecke{3cm} & el motivo & \luecke{3cm} \\[0.5em]
\hline
el trabajo & \luecke{3cm} & el plan & \luecke{3cm} \\[0.5em]
\hline
\end{tabularx}

\vspace{0.8em}

\aufgabe{Wortfamilien: Finde die verwandten Wörter.}{6}

\begin{tabularx}{\textwidth}{@{}|X|X|X|@{}}
\hline
\rowcolor{hellgrau}
\textbf{Verb} & \textbf{Substantiv (-ción)} & \textbf{Adjektiv} \\
\hline
educar & \luecke{3cm} & educativo \\[0.5em]
\hline
\luecke{3cm} & participación & participativo \\[0.5em]
\hline
administrar & \luecke{3cm} & \luecke{3cm} \\[0.5em]
\hline
\end{tabularx}

\vspace{0.8em}

\aufgabe{Diskussion: ¿Qué hace un buen jefe? (Buch S. 104, Aufg. 5)}{8}

\begin{buchverweis}
\textbf{Notiere 5 Eigenschaften} eines guten Chefs. Verwende die neu gelernten Substantive!
\end{buchverweis}

\vspace{0.3em}
\textbf{Un buen jefe necesita...}

\begin{enumerate}[leftmargin=2em]
    \item \luecke{10cm}
    \item \luecke{10cm}
    \item \luecke{10cm}
    \item \luecke{10cm}
    \item \luecke{10cm}
\end{enumerate}

\vspace{0.5em}
\textbf{Begründe eine Eigenschaft mit \es{porque} oder \es{ya que}:}

\vspace{0.3em}
\begin{tabularx}{\textwidth}{@{}|X|@{}}
\hline
\\[0.4em]
\luecke{14cm} \\[0.7em]
\luecke{14cm} \\[0.4em]
\hline
\end{tabularx}

\vspace{1em}
\hrule
\vspace{0.3em}
\begin{flushright}
\textbf{Gesamt: \luecke{1cm} / 32 Punkte}
\end{flushright}

\end{document}
